\documentclass[12pt]{article}
\usepackage[utf8]{inputenc}
\usepackage[T1]{fontenc}
\usepackage{hyperref}
\title{Mayi-Kulan Language Resources for K-12 Curriculum Development}
\author{Compiled by Claude Sonnet 5.5 (claude-sonnet-5-5) for Andr\'as Kornai}
\date{}
\begin{document}
\maketitle

\section{Introduction}
Mayi-Kulan (ISO 639-3: xyk) is an extinct language of the Mayi group of the Queensland Gulf Country. Sources disagree on its location. Wikipedia places it on Cape York Peninsula, but AIATSIS and Breen's work place the Mayi languages in the Gulf Country. Breen treats Mayi-Kulan as one dialect in a chain that includes Ngawun, Mayi-Yapi, Mayi-Thakurti and Wunumara. He describes the whole group as very sparsely recorded, with no competent speakers for many years.

\section{Primary Educational Resources}
No primers, textbooks or school programmes were found.

\section{Linguistic and Reference Materials}
\begin{itemize}
\item Gavan Breen (1981), \emph{The Mayi Languages of the Queensland Gulf Country} (Glottolog reference record). \url{https://glottolog.org/resource/reference/id/64840}
\item Gavan Breen, review of Laves's Mayi-Thakurti material (PDF, ANU). \url{https://www0.anu.edu.au/linguistics/nash/aust/laves/Breen-MT.pdf}
\item Wikipedia, Mayi-Kulan language. It has a short word list drawn from Blake (1981) and should be treated with caution given the location discrepancy above. \url{https://en.wikipedia.org/wiki/Mayi-Kulan_language}
\item Glottolog record for the ISO code mnt, which groups the Mayi varieties. \url{https://glottolog.org/resource/languoid/iso/mnt}
\end{itemize}

\section{Religious and Community Resources}
No community or religious resources specific to Mayi-Kulan were found.

\section{Digital and Audio Resources}
No audio, video or software resources were found.

\section{Curriculum Development Context}
Very little is documented, so a curriculum could not teach the language itself in any substantial way. It could use the small word list and Breen's comparative material. It could also teach about the Mayi languages and the Gulf Country, with descendant communities leading and any archival access arranged through AIATSIS.
\end{document}